\documentclass[12pt, a4paper]{article}

\usepackage{fontspec}
\usepackage{polyglossia}
\setdefaultlanguage{russian}
\setotherlanguage{english}
\setmainfont{Times New Roman}
\setmonofont{Consolas}
\newfontfamily\cyrillicfonttt{Consolas}

\usepackage[left=2.5cm, right=1.5cm, top=2cm, bottom=2cm]{geometry}
\usepackage{setspace}\onehalfspacing
\usepackage{xcolor}
\usepackage[most]{tcolorbox}
\usepackage{hyperref}
\usepackage{enumitem}
\usepackage{booktabs}
\usepackage{indentfirst}

\hypersetup{
    colorlinks=true,
    linkcolor=blue!60!black,
    urlcolor=blue!60!black
}

\newtcolorbox{promptbox}[1][]{
    colback=blue!3,
    colframe=blue!40!black,
    fonttitle=\bfseries,
    title=Текст промпта,
    #1
}

\newtcolorbox{responsebox}[1][]{
    colback=green!3,
    colframe=green!35!black,
    fonttitle=\bfseries,
    title=Ответ LLM,
    #1
}

\newtcolorbox{commentbox}[1][]{
    colback=orange!5,
    colframe=orange!40!black,
    fonttitle=\bfseries,
    title=Мой комментарий,
    #1
}

\begin{document}


\section{Введение}

\textbf{Датасет:} из соревнования Diabetes Prediction Challenge (Kaggle)\\
\textbf{Источник:} \url{https://www.kaggle.com/competitions/playground-series-s5e12/data}

Датасет содержит медицинские и демографические данные и предназначен для бинарной классификации --- предсказания наличия диагностированного диабета \

(\texttt{diagnosed\_diabetes}: 0 или 1). 

В наборе 26 столбцов: числовые показатели здоровья (ИМТ, давление, пульс, холестерин), 
факторы образа жизни (физическая активность, диета, сон, экранное время), 
демографические признаки (пол, этничность, образование, доход) и медицинский анамнез 
(семейная история диабета, гипертензия, сердечно-сосудистые заболевания).

\textbf{LLM:} Claude (Anthropic).

\section{Промпт 1: Понимание структуры данных}

\begin{promptbox}
У меня есть датасет для предсказания диабета (train.csv) с 700\,000 строк и 26 столбцами. Столбцы: id, age, alcohol\_consumption\_per\_week, physical\_activity\_minutes\_per\_week, diet\_score, sleep\_hours\_per\_day, screen\_time\_hours\_per\_day, bmi, waist\_to\_hip\_ratio, systolic\_bp, diastolic\_bp, heart\_rate, cholesterol\_total, hdl\_cholesterol, ldl\_cholesterol, triglycerides, gender, ethnicity, education\_level, income\_level, smoking\_status, employment\_status, family\_history\_diabetes, hypertension\_history, cardiovascular\_history, diagnosed\_diabetes. Целевая переменная --- diagnosed\_diabetes (0/1). Опиши структуру данных: какие типы признаков присутствуют, какие из них числовые, какие категориальные, какие бинарные. Предположи, какие проблемы с качеством данных могут встретиться.
\end{promptbox}

\begin{responsebox}
LLM классифицировала 26 столбцов следующим образом:

\textbf{Числовые непрерывные (15):} age, alcohol\_consumption\_per\_week, physical\_activity\_minutes\_per\_week, diet\_score, sleep\_hours\_per\_day, screen\_time\_hours\_per\_day, bmi, waist\_to\_hip\_ratio, systolic\_bp, diastolic\_bp, heart\_rate, cholesterol\_total, hdl\_cholesterol, ldl\_cholesterol, triglycerides.

\textbf{Категориальные номинальные (6):} gender (Male/Female/Other), ethnicity (5 значений), education\_level (4 значения), income\_level (5 значений), smoking\_status (3 значения), employment\_status (4 значения).

\textbf{Бинарные (4):} family\_history\_diabetes, hypertension\_history, cardiovascular\_history, diagnosed\_diabetes.

\textbf{Технический (1):} id.

LLM предположила возможные проблемы с качеством: пропущенные значения, выбросы в BMI и артериальном давлении, дисбаланс классов в целевой переменной, мультиколлинеарность между компонентами холестерина (cholesterol\_total $\approx$ HDL + LDL + triglycerides/5 по формуле Фридвальда), необходимость кодирования категориальных признаков.
\end{responsebox}

\begin{commentbox}
LLM верно предположила мультиколлинеарность компонентов холестерина по формуле Фридвальда, что действительно может привести к нестабильности коэффициентов в линейных моделях. Однако модель ошиблась в предположении о пропущенных значениях --- при проверке данных выяснилось, что в датасете нет ни одного пропуска, что нетипично для реальных медицинских данных и указывает на предварительную очистку или синтетическую генерацию. Также LLM не обратила внимания на непропорционально большой объём данных (700\,000 записей), нехарактерный для клинических исследований, что дополнительно подтверждает искусственную природу набора.
\end{commentbox}

\section{Промпт 2: Тренды образа жизни и диабет}

\begin{promptbox}
Проанализируй связь между образом жизни и диагнозом диабета в моём датасете. Вот статистика: для группы с диабетом (diagnosed\_diabetes=1, n=436\,307): средняя физическая активность 73.5 мин/нед, средний diet\_score 5.91, средний BMI 26.1, среднее потребление алкоголя 2.07 ед/нед, средний сон 7.00 ч/день, среднее экранное время 6.04 ч/день. Для группы без диабета (n=263\,693): средняя физическая активность 91.4 мин/нед, средний diet\_score 6.06, средний BMI 25.5, алкоголь 2.07, сон 7.00, экранное время 5.97. Какие тренды ты видишь? Какие факторы наиболее выражены?
\end{promptbox}

\begin{responsebox}
LLM выделила следующие тренды:

\begin{enumerate}
    \item \textbf{Физическая активность} --- сильнейший дифференциатор: группа с диабетом в среднем тратит 73.5 мин/нед против 91.4 мин/нед у здоровых, что составляет снижение на 24\%. Это согласуется с медицинской литературой о сидячем образе жизни как факторе риска.
    \item \textbf{BMI} показывает умеренную, но устойчивую разницу (26.1 vs 25.5) --- оба значения вблизи порога избыточного веса (25), но группа с диабетом его пересекает.
    \item \textbf{Diet score} различается минимально (5.91 vs 6.06), что говорит о слабой дискриминативной способности качества диеты.
    \item \textbf{Потребление алкоголя} идентично (2.07 vs 2.07) --- тренд отсутствует.
    \item \textbf{Сон и экранное время} практически не различаются между группами.
\end{enumerate}

Также модель отметила, что общая распространённость диабета составляет 62.3\%, что крайне нереалистично для генеральной популяции.
\end{responsebox}

\begin{commentbox}
Наиболее значимым наблюдением является разница в физической активности: 24\% снижение у группы с диабетом --- это клинически значимый результат, согласующийся с рекомендациями ВОЗ о минимуме 150 минут умеренной активности в неделю (обе группы далеки от этого порога). LLM справедливо указала на подозрительно высокую распространённость диабета в 62.3\%, тогда как в реальных популяциях этот показатель составляет 8--12\% по данным ВОЗ, что подтверждает гипотезу о синтетической генерации с искусственно смещённым балансом классов. Удивительно, что потребление алкоголя абсолютно не различается между группами --- в реальности связь, как правило, нелинейная (U-образная кривая), но при усреднении эффект теряется.
\end{commentbox}

\section{Промпт 3: Обнаружение аномалий в данных о здоровье}

\begin{promptbox}
В моём датасете о диабете обнаружены следующие диапазоны значений: physical\_activity\_minutes\_per\_week: min=1, max=747, mean=80.2; BMI: min=15.1, max=38.4; systolic\_bp: min=91, max=163; heart\_rate: min=42, max=101; sleep\_hours\_per\_day: min=3.1, max=9.9; screen\_time\_hours\_per\_day: min=0.6, max=16.5; waist\_to\_hip\_ratio: min=0.68, max=1.05. Также: 7\,244 записи с активностью >300 мин/нед, 522 записи с BMI>35, 41 запись с экранным временем >14 ч/день. Выяви аномалии и потенциальные выбросы. Какие значения выглядят подозрительно с медицинской точки зрения?
\end{promptbox}

\begin{responsebox}
LLM проанализировала диапазоны с медицинской точки зрения:

\begin{enumerate}
    \item \textbf{Физическая активность max=747 мин/нед} ($\approx$107 мин/день, 7 дней) --- экстремально, но физически возможно для профессиональных спортсменов. 7\,244 записи с >300 мин/нед (>43 мин/день) вполне правдоподобны для активных людей.
    \item \textbf{Экранное время max=16.5 ч/день} --- оставляет лишь 7.5 часов на сон и другие активности; подозрительно, но технически возможно.
    \item \textbf{BMI min=15.1} --- соответствует тяжёлому дефициту массы тела (граница анорексии); max=38.4 --- ожирение II степени, но подозрительно <<обрезан>> сверху (в реальных данных встречаются значения до 50+).
    \item \textbf{Пульс 42--101 уд/мин} --- 42 уд/мин характерен для брадикардии (норма у спортсменов), 101 --- пограничная тахикардия. Диапазон подозрительно узкий для клинических данных.
    \item \textbf{Систолическое АД 91--163 мм рт. ст.} --- 91 --- пограничная гипотензия, 163 --- гипертензия II стадии. Отсутствие экстремальных значений нетипично.
    \item \textbf{Сон 3.1--9.9 часов} --- ни одного значения ниже 3 или выше 10, что указывает на искусственное ограничение диапазона.
\end{enumerate}
\end{responsebox}

\begin{commentbox}
LLM продемонстрировала сильную сторону --- способность сопоставить числовые диапазоны с медицинскими нормами, верно отметив, что пульс 42 уд/мин характерен для тренированных спортсменов, а ИМТ 15.1 соответствует тяжёлому дефициту массы тела. Ключевое наблюдение, которое LLM могла бы подчеркнуть сильнее: все числовые признаки имеют подозрительно <<чистые>> диапазоны без экстремальных выбросов, что нетипично для реальных медицинских данных, где встречаются ошибки ввода и артефакты измерений. Ограничение BMI значением 38.4 сверху особенно показательно --- в клинической практике значения выше 40 (морбидное ожирение) достаточно распространены, а их отсутствие подтверждает генерацию данных по заранее заданным распределениям.
\end{commentbox}

\section{Промпт 4: Демографические паттерны и медицинский анамнез}

\begin{promptbox}
Вот данные о распространённости диабета по подгруппам в моём датасете (700\,000 записей). По семейному анамнезу: с семейной историей диабета --- 86.7\% (n=104\,581), без --- 58.0\% (n=595\,419). По гипертензии: с историей --- 65.4\% (n=127\,393), без --- 61.6\%. По сердечно-сосудистым заболеваниям: с историей --- 69.1\% (n=21\,227), без --- 62.1\%. По полу: женщины 62.2\%, мужчины 62.4\%, другое 64.1\%. По этничности: White 62.4\%, Hispanic 61.6\%, Black 62.4\%, Asian 62.8\%, Other 63.6\%. По курению: Current 62.3\%, Never 62.3\%, Former 62.5\%. Проанализируй, какие факторы сильнее всего влияют на диабет. Сформулируй практические выводы для системы здравоохранения.
\end{promptbox}

\begin{responsebox}
LLM провела анализ:

\begin{enumerate}
    \item \textbf{Семейная история --- доминирующий предиктор}: разница в 28.7 п.п. (86.7\% vs 58.0\%). Это сильнейший одиночный фактор в датасете.
    \item \textbf{Сердечно-сосудистые заболевания} показывают значимую разницу (69.1\% vs 62.1\%), как и \textbf{гипертензия} (65.4\% vs 61.6\%).
    \item \textbf{Пол, этничность и статус курения практически не различаются} (разброс 1--2 п.п.), что означает отсутствие дискриминативной способности.
\end{enumerate}

\textbf{Практические рекомендации:}
\begin{itemize}
    \item Таргетный скрининг для лиц с семейной историей диабета.
    \item Интегрированная оценка метаболического и кардиоваскулярного риска.
    \item Универсальный скрининг вне зависимости от пола/этничности, поскольку риск распределён равномерно.
\end{itemize}
\end{responsebox}

\begin{commentbox}
Данные выявляют парадоксальную картину: семейная история диабета является мощнейшим предиктором (разница 28.7 п.п.), тогда как традиционно значимые факторы --- пол, этничность, курение --- практически не дифференцируют группы. В реальной эпидемиологии этнические различия в распространённости диабета II типа хорошо документированы (например, повышенный риск у афроамериканцев и латиноамериканцев по данным CDC), поэтому их отсутствие является ещё одним признаком синтетической генерации данных. LLM сформулировала разумные рекомендации о таргетном скрининге, но не смогла критически оценить нереалистичность равномерного распределения риска по демографическим группам, принимая данные за истину без дополнительных оговорок.
\end{commentbox}

\section{Промпт 5: Итоговые выводы и рекомендации по моделированию}

\begin{promptbox}
На основании всего анализа датасета предсказания диабета (700\,000 строк, 26 признаков), резюмируй ключевые находки и предложи: (1) какие 5 признаков наиболее важны для предсказания диабета и почему; (2) какие признаки можно исключить без потери качества; (3) какие шаги предобработки необходимы; (4) какой тип модели рекомендуешь и почему. Учитывай дисбаланс классов 62.3\% к 37.7\%.
\end{promptbox}

\begin{responsebox}
LLM предложила:

\textbf{Топ-5 признаков:}
\begin{enumerate}
    \item \texttt{family\_history\_diabetes} --- сильнейший одиночный предиктор (86.7\% vs 58.0\%).
    \item \texttt{physical\_activity\_minutes\_per\_week} --- 24\% разница между группами.
    \item \texttt{bmi} --- клинически значимая разница (26.1 vs 25.5).
    \item \texttt{age} --- разница в 4 года между средними.
    \item \texttt{hypertension\_history} --- сигнал коморбидности.
\end{enumerate}

\textbf{Признаки для исключения:} id (технический), один из коллинеарных компонентов холестерина (оставить HDL и триглицериды, исключить cholesterol\_total и LDL), smoking\_status (нет дискриминативной способности), gender и ethnicity (разница $<$2 п.п.).

\textbf{Предобработка:} one-hot кодирование категориальных признаков, стандартизация числовых, создание интерактивных признаков (age $\times$ bmi, family\_history $\times$ physical\_activity).

\textbf{Рекомендация по модели:} Gradient Boosting (XGBoost/LightGBM) как основная модель; логистическая регрессия как интерпретируемый бейзлайн. При дисбалансе 62/38 достаточно стратифицированных разбиений и метрик F1/AUC-ROC, агрессивный ресемплинг не нужен.
\end{responsebox}

\begin{commentbox}
LLM предложила практически применимую стратегию моделирования, однако допустила характерную ошибку --- рекомендуя удалить \texttt{cholesterol\_total} из-за мультиколлинеарности, она не учла, что для ансамблевых моделей (XGBoost) мультиколлинеарность не является проблемой, в отличие от линейной регрессии. Это типичный пример переноса знаний из одного контекста в другой без адаптации. Рекомендация по инженерии признаков (age $\times$ bmi, family\_history $\times$ physical\_activity) отражает понимание медицинского контекста, где взаимодействие факторов часто важнее их индивидуального вклада. Стоит отметить, что LLM верно определила умеренный дисбаланс классов и не стала рекомендовать агрессивные методы ресемплинга (SMOTE), что свидетельствует о практичном подходе.
\end{commentbox}


\section{Заключение}

В ходе работы было проведено разведочное исследование датасета предсказания диабета (700\,000 записей, 26 признаков) с помощью пяти целенаправленных промптов, охвативших: понимание структуры данных, анализ трендов образа жизни, обнаружение аномалий, демографический анализ и выработку рекомендаций по моделированию.

LLM оказалась наиболее полезна для быстрой генерации гипотез и интерпретации данных в медицинском контексте --- модель корректно идентифицировала пороги брадикардии, формулу Фридвальда для холестерина, рекомендации ВОЗ по физической активности.

Ключевое ограничение: LLM не способна самостоятельно вычислять статистику --- она полностью полагается на предоставленные ей числа и может уверенно оперировать неверными данными, если ей их передать. Кроме того, модель склонна принимать данные за <<чистую монету>>, не подвергая сомнению их происхождение: LLM не инициативно указала на нереалистичную распространённость диабета в 62.3\% и на отсутствие этнических различий, которые хорошо документированы в реальной эпидемиологии.

Оптимальный рабочий процесс для EDA --- комбинация LLM для генерации гипотез и интерпретации с Python/Pandas для фактических вычислений и проверки. LLM ускоряет анализ, но не заменяет его --- каждый вывод модели требует верификации на реальных данных.

\end{document}
